

%------------------------------------------------------------
\section{Index bitmap}
%------------------------------------------------------------
\begin{frame}{Le problème}

Comment indexer une table sur un attribut
qui ne prend qu'un petit nombre de valeurs\,?
\begin{itemize}
  \item Avec un arbre B\,: pas très bon car chaque valeur
      est peu sélective
 \item Avec un hachage\,: pas très bon non plus
          car il y a beaucoup de collisions.
\end{itemize}

Or situation fréquente, notamment dans les entrepôts de données

\end{frame}
%------------------------------------------------------------
\begin{frame}{Exemple\,: codification des films}

\begin{center}
\begin{tabular}{l|l|c|c}
rang & titre & genre  & ...  \\
\hline
1 & Vertigo  & Suspense & ... \\
2 & Brazil & Science-Fiction & ... \\
3 & Twin Peaks & Fantastique & ... \\
4 & Underground & Drame & ... \\
5 & Easy Rider &  Drame & ...\\
6 & Psychose & Drame & ... \\
7 & Greystoke &  Aventures & ... \\
8 & Shining & Fantastique & ... \\
... & ... &  ... & ...\\
\hline
\end{tabular}
\end{center}

\end{frame}
%------------------------------------------------------------
\begin{frame}{Principes de l'index bitmap}

Soit un attribut $A$, prenant $n$ valeurs possibles $[v_1, \ldots v_n]$
\begin{itemize}
  \item On crée $n$ tableaux de bit, un pour chaque valeur $v_i$ 

   \item Ce tableau contient un bit pour  chaque enregistrement $e$

   \item   Le bit d'un enregistrement $e$ est à 1 si $e.A=v_i$,
                à 0 sinon
\end{itemize}

\end{frame}
%------------------------------------------------------------
\begin{frame}{Exemple}

\begin{center}
\begin{tabular}{|l|c|c|c|c|c|c|c|c|}
\hline
 & 1 & 2 & 3 & 4 & 5 & 6 & 7 & 8\\
\hline
Drame & 0 & 0 & 0 & 1 & 1 & 1 & 0 & 0 \\
Science-Fiction  & 0 & 1 & 0 & 0 & 0 & 0 & 0 & 0 \\
Comédie  & 0 & 0 & 0 & 0 & 0 & 0 & 0 & 0 \\
\hline
\end{tabular}

\begin{tabular}{|l|c|c|c|c|c|c|c|c|}
\hline
 & 9 & 10 & 11 & 12 & 13 & 14 & 15 & 16\\
\hline
Drame &  0 & 0 & 0 & 0 & 0 & 0 & 1 & 0\\
Science-Fiction  & 0 & 1 & 1 & 0 & 0 & 0 & 0 & 0\\
Comédie  &  1 & 0 & 0 & 1 & 0 & 0 & 0 & 1\\
\hline
\end{tabular}
\end{center}

\end{frame}
%------------------------------------------------------------
\begin{frame}{Recherche}

Soit une requête comme\,:\\
\texttt{SELECT * FROM Film WHERE genre='Drame'}

\begin{itemize}
  \item On prend le tableau pour la valeur \texttt{Drame}

   \item On garde toutes les cellules à 1

   \item On accède aux enregistrements par l'adresse
\end{itemize}

=$>$ très efficace si $n$, le nombre de valeurs, est petit.

\end{frame}
%------------------------------------------------------------
\begin{frame}{Autre exemple}

\texttt{SELECT COUNT(*) FROM   Film
WHERE  genre IN ('Drame', 'Comédie')}

\begin{itemize}
  \item On compte le nombre de 1 dans le tableau \texttt{Drame}
  \item On compte le nombre de 1 dans le tableau \texttt{Comédie}
   \item On fait la somme et c'est fini
\end{itemize}

=$>$ typique des requêtes \textit{DataWarehouse}

\end{frame}
%------------------------------------------------------------
\begin{frame}{Indexation de données géométriques}

Certaines applications stockent des données géométriques
(ex, cadastre, réseaux routiers, etc.). Requêtes
courantes\,:
\begin{enumerate}
  \item pointé,
   \item fenêtrage,
  \item jointure spatiale. 
\end{enumerate}
Comment optimiser ces opérations avec des index\,?

\end{frame}
%------------------------------------------------------------
\begin{frame}{Obstacles}

Principaux obstacles à l'utilisation d'un index standard\,:
\begin{enumerate}
  \item On ne peut pas utiliser un arbre-B

      car 
      pas d'ordre sur les points ou les polygones.

  \item Opérations coûteuses

      Algorithmes géométriques, souvent de complexité
            polynomiale, difficiles à implanter.
\end{enumerate}

$\Rightarrow$ des structures spécialisées pour les
données géométriques

\end{frame}
%------------------------------------------------------------
\begin{frame}{Principes}

Pour limiter le coût des opérations, on 
utilise une approximation rectangulaire 
des données géométriques\,: le rectangle englobant.

\begin{enumerate}
  \item l'index est construit sur l'ensemble des rectangles\,;
  \item On effectue un {\em filtrage} basé sur les rectangles
         pour les pointés, fenêtrages et jointures.
  \item Quand on accède aux objets eux-mêmes, on effectue
          une seconde opération sur la vraie géométrie
          (étape de {\em raffinement})
\end{enumerate}


\end{frame}
%------------------------------------------------------------
\begin{frame}{L'arbre R, définition}

Assez proche de l'arbre-B dans ses principes\,:

\begin{enumerate}
  \item c'est un arbre équilibré\,;
  \item chaque entrée est une paire $[adresse, rectangle]$\,;
  \item tous les n{\oe}uds sont occupés au moins à 50\,\%
  \item le rectangle d'une entrée  $[adresse, rectangle]$
          englobe {\em tous} les rectangles du n{\oe}ud
           pointé par $adresse$\,;
\end{enumerate}


\end{frame}
%------------------------------------------------------------
\begin{frame}{L'arbre R, exemple}


\figSlideWithSize{rtree}{11cm}


\end{frame}
%------------------------------------------------------------
\begin{frame}{L'arbre R, pointé}


\figSlideWithSize{pqRtree}{11cm}


\end{frame}
%------------------------------------------------------------
\begin{frame}{L'arbre R, insertion (1)}


\figSlideWithSize{insRtree1}{11cm}


\end{frame}
%------------------------------------------------------------
\begin{frame}{L'arbre R, insertion (2)}


\figSlideWithSize{insRtree2}{11cm}


\end{frame}
